% defer/rcufundamental.tex
% mainfile: ../perfbook.tex
% SPDX-License-Identifier: CC-BY-SA-3.0

\subsection{RCU基础}
\label{sec:defer:RCU Fundamentals}
\OriginallyPublished{sec:defer:RCU Fundamentals}{RCU Fundamentals}{Linux Weekly News}{PaulEMcKenney2007WhatIsRCUFundamentally}

本节重新审视前一节涵盖的内容，但不依赖于任何特定的示例或用例。
喜欢贴近实际代码的读者可以跳过本节所介绍的底层基础知识。

RCU保护链式数据结构的常见用法由三个基本机制组成，第一个用于插入，
第二个用于删除，第三个用于允许读者容忍并发的插入和删除操作。
\Cref{sec:defer:Publish-Subscribe Mechanism}
描述了用于插入的publish-subscribe机制，
\cref{sec:defer:Wait For Pre-Existing RCU Readers}
描述了等待已有RCU读者如何实现删除，\footnote{
	是的，RCU就是这样，这个最基本的RCU组件不是第一个介绍，
	而是第二个。}
而
\cref{sec:defer:Maintain Multiple Versions of Recently Updated Objects}
讨论了维护最近更新对象的多个版本如何允许并发的插入和删除操作。
最后，
\cref{sec:defer:Summary of RCU Fundamentals}
总结了RCU的基础知识。

\subsubsection{发布-订阅机制}
\label{sec:defer:Publish-Subscribe Mechanism}

由于RCU读者不会被RCU更新者排斥，RCU保护的数据结构可能在读者
访问时发生变化。
被访问的数据项可能被移动、移除或替换。
因为数据结构不会为读者"保持静止"，所以每个读者的访问可以被视为
订阅（subscribe）RCU保护数据项的当前版本。
就更新者而言，它们可以被视为发布（publish）新版本。

% @@@ Merge usage section into "Which to Choose?" and suggest choices.

\begin{figure}
\centering
\resizebox{3in}{!}{\includegraphics{defer/pubsub}}
\caption{发布/订阅约束}
\label{fig:defer:Publication/Subscription Constraints}
\end{figure}

不幸的是，正如
\cref{sec:toolsoftrade:Shared-Variable Shenanigans}
中所述以及
\cref{sec:defer:Minimal Insertion and Deletion}
中重申的那样，
对这些发布和订阅操作使用\IXplx{plain access}{es}是不明智的。
相反，必须告知编译器和CPU需要小心处理，
如\cref{fig:defer:Publication/Subscription Constraints}
所示，该图说明了
\cref{lst:defer:Insertion and Deletion With Concurrent Readers}
中\co{ins_route()}（及其调用者）和\co{access_route()}并发执行之间的交互。

\cref{fig:defer:Publication/Subscription Constraints}
中的\co{ins_route()}列显示了\co{ins_route()}的调用者分配一个新的
\co{route}结构，此时该结构包含的是初始化前的垃圾数据。
调用者随后初始化新分配的结构，然后调用\co{ins_route()}发布指向
新\co{route}结构的指针。
发布不会影响结构的内容，因此结构内容在发布后仍然有效。

同一图中的\co{access_route()}列显示了指针被订阅和解引用的过程。
这个解引用操作绝对必须看到一个有效的\co{route}结构，而不是初始化
前的垃圾数据，因为引用垃圾数据可能导致内存损坏、崩溃和挂起。
如前所述，避免此类垃圾数据意味着发布和订阅操作必须告知编译器和CPU
维持所需排序的必要性。

发布通过\co{rcu_assign_pointer()}执行，它确保\co{ins_route()}调用者的
初始化操作排在实际发布操作存储指针之前。
此外，\co{rcu_assign_pointer()}必须是原子的，即并发读者看到的要么是
指针的旧值，要么是新值，而不是两者的某种混合。
这些要求由C11 store-release操作满足，实际上在Linux内核中，
\co{rcu_assign_pointer()}是用\co{smp_store_release()}定义的，
后者类似于C11 store-release。

注意，如果需要并发更新，则需要某种同步机制来协调同一指针上的
多个并发\co{rcu_assign_pointer()}调用。
在Linux内核中，锁是首选机制，但几乎任何同步机制都可以使用。
一个特别轻量级的同步机制的例子是
\cref{chp:Data Ownership}的数据所有权（data ownership）：
如果每个指针由特定线程拥有，则该线程可以在没有额外同步开销的情况下
对该指针执行\co{rcu_assign_pointer()}。

\QuickQuiz{
	对RCU更新者使用数据所有权是否意味着更新可以使用与
	相应单线程代码完全相同的指令序列？
}\QuickQuizAnswer{
	有时候是的，例如在TSO系统（如x86或IBM大型机）上，
	store-release操作只发出一条存储指令。
	然而，弱排序系统还必须发出memory barrier或者
	store-release指令。
	此外，移除数据需要相当多的额外工作，因为在释放被移除的数据
	之前必须等待已有的读者完成。
}\QuickQuizEnd

订阅通过\co{rcu_dereference()}执行，它防止编译器和（在一种情况下）
CPU将指针加载与后续解引用该指针的内存操作重排序。
与\co{rcu_assign_pointer()}类似，\co{rcu_dereference()}必须是原子的，
即加载的值必须来自单次存储，例如，编译器不能拆分该加载。\footnote{
	也就是说，编译器不能将该加载拆分为多个更小的加载，
	如\cref{sec:toolsoftrade:Shared-Variable Shenanigans}中
	"load tearing"下所述。}
不幸的是，编译器对\co{rcu_dereference()}的支持充其量还在进行中~\cite{PaulEMcKennneyConsumeP0190R4,PaulEMcKenney2017markconsumeP0462R1,JFBastien2018P0750R1consume}。
在此期间，Linux内核依赖volatile加载、各种CPU架构的细节、
编码限制~\cite{PaulEMcKenney2014rcu-dereference}，
以及在DEC Alpha~\cite{ALPHA2002}上的一条memory-barrier指令。
然而，在其他架构上，\co{rcu_dereference()}通常只发出一条加载指令，
与等效的单线程代码相同。
编码限制在
\cref{sec:memorder:Address- and Data-Dependency Difficulties}
中有更详细的描述，不过，常见的字段选择（\qtco{->}）情况工作得很好。
不需要极致读侧性能的软件可以改用C11 \IXpl{acquire load}，
它提供了所需的排序及更多保证，但需要付出一定代价。
希望更轻量级的\co{rcu_dereference()}编译器支持将在适当时候出现。

简而言之，使用\co{rcu_assign_pointer()}发布指针和使用
\co{rcu_dereference()}订阅指针成功避免了
\cref{fig:defer:Publication/Subscription Constraints}
中描述的"Not OK"垃圾加载。
因此，这两个原语可以用于向链式结构添加新数据，而不会干扰并发读者。

\QuickQuiz{
	但假设更新者正在向链表中添加和移除多个数据项，
	同时一个读者正在遍历同一链表。
	具体来说，假设链表最初包含元素A、B和~C，
	更新者移除了元素A，然后在链表末尾添加了新元素D。
	读者很可能看到\{A, B, C, D\}，而这个元素序列实际上
	从未存在过！
	在哪个平行宇宙中这算得上"不干扰并发读者"？？？
}\QuickQuizAnswer{
	在这样一个世界中：遍历读者只需要遍历在整个遍历期间
	一直存在的元素。
	在本例中，那就是元素B和~C\@。
	因为元素A和~D各自只在遍历的一部分时间内存在，读者可以
	遍历它们，但并非必须遍历。
	注意，这支持了常见的情况，即读者只是查找单个项目，
	并不知道或关心其他项目的存在与否。

	如果需要更强的一致性，则需要更高代价的同步机制，
	例如sequence locking或读写锁。
	但如果不需要更强的一致性（实际上很多时候不需要），
	为什么要付出更高的代价呢？
}\QuickQuizEnd

在不干扰读者的情况下向链式结构添加数据是一件好事，
与单线程读者相比可以不增加读侧开销的情况更是如此。
然而，在大多数情况下还需要移除数据，这是下一节的主题。

\subsubsection{等待已有的RCU读者}
\label{sec:defer:Wait For Pre-Existing RCU Readers}

RCU最基本的组件是等待事情完成。
当然，还有很多其他方式来等待事情完成，包括reference counting、
读写锁、事件等。
RCU的巨大优势在于它可以等待（比如说）20,000个不同的事情完成，
而不必显式跟踪每一个，也不必担心显式跟踪方案固有的性能下降、
可扩展性限制、复杂的死锁场景和内存泄漏风险。

在RCU的情况下，等待的每一件事称为
\emph{\IXBhmr{RCU read-side}{critical section}}。
如\cref{tab:defer:Core RCU API}中所述，
RCU read-side critical section以\co{rcu_read_lock()}原语开始，
以相应的\co{rcu_read_unlock()}原语结束。
RCU read-side critical section可以嵌套，并且可以包含几乎任何代码，
只要该代码不包含quiescent state。
例如，在Linux内核中，在RCU read-side critical section中睡眠是
不允许的，因为上下文切换是一个quiescent state。\footnote{
	然而，一种称为SRCU~\cite{PaulEMcKenney2006c}的特殊RCU形式
	确实允许在SRCU read-side critical section中进行一般性睡眠。}
如果你遵守这些约定，你可以使用RCU来等待\emph{任何}已有的
RCU read-side critical section完成，而
\co{synchronize_rcu()}使用间接方式来执行实际的
等待~\cite{MathieuDesnoyers2012URCU,McKenney:2013:SDS:2483852.2483867}。

\begin{figure}
\centering
\resizebox{3in}{!}{\includegraphics{defer/RCUGuaranteeFwd}}
\caption{RCU读者与后续Grace Period}
\label{fig:defer:RCU Reader and Later Grace Period}
\end{figure}

RCU read-side critical section与后续RCU grace period之间的关系
是一种"如果-那么"关系，如
\cref{fig:defer:RCU Reader and Later Grace Period}所示。
如果给定critical section的任何部分先于给定grace period的开始，
那么RCU保证该critical section的全部将先于该grace period的结束。
在图中，\co{P0()}对\co{x}的访问先于\co{P1()}对同一变量的访问，
因此也先于\co{P1()}调用\co{synchronize_rcu()}产生的grace period。
因此可以保证\co{P0()}对\co{y}的访问将先于\co{P1()}的访问。
在这种情况下，如果\co{r1}的最终值为0，则\co{r2}的最终值
也保证为0。

\QuickQuiz{
	在\cref{fig:defer:RCU Reader and Later Grace Period}中，
	\co{r1}和\co{r2}还有哪些其他可能的最终值？
}\QuickQuizAnswer{
	\co{r1 == 0 && r2 == 0}的可能性已在正文中指出。
	鉴于\co{r1 == 0}意味着\co{r2 == 0}，我们知道
	\co{r1 == 0 && r2 == 1}是被禁止的。
	以下讨论将表明\co{r1 == 1 && r2 == 1}和
	\co{r1 == 1 && r2 == 0}都是可能的。
}\QuickQuizEnd

\begin{figure}
\centering
\resizebox{3in}{!}{\includegraphics{defer/RCUGuaranteeRev}}
\caption{RCU读者与更早的Grace Period}
\label{fig:defer:RCU Reader and Earlier Grace Period}
\end{figure}

RCU read-side critical section与更早的RCU grace period之间的关系
也是一种"如果-那么"关系，如
\cref{fig:defer:RCU Reader and Earlier Grace Period}所示。
如果给定critical section的任何部分后于给定grace period的结束，
那么RCU保证该critical section的全部将后于该grace period的开始。
在图中，\co{P0()}对\co{y}的访问后于\co{P1()}对同一变量的访问，
因此后于\co{P1()}调用\co{synchronize_rcu()}产生的grace period。
因此可以保证\co{P0()}对\co{x}的访问将后于\co{P1()}的访问。
在这种情况下，如果\co{r2}的最终值为1，则\co{r1}的最终值
也保证为1。

\QuickQuiz{
	如果\co{P0()}在
	\cref{fig:defer:RCU Reader and Earlier Grace Period}
	中的两次访问顺序颠倒，会发生什么？
}\QuickQuizAnswer{
	绝对不会有任何变化。
	\co{P0()}对\co{x}和\co{y}的加载在同一个
	RCU read-side critical section中这一事实就足够了；
	它们的顺序无关紧要。
}\QuickQuizEnd

\begin{figure}
\centering
\resizebox{3in}{!}{\includegraphics{defer/RCUGuaranteeMid}}
\caption{Grace Period内的RCU读者}
\label{fig:defer:RCU Reader Within Grace Period}
\end{figure}

最后，如\cref{fig:defer:RCU Reader Within Grace Period}所示，
RCU read-side critical section可以被RCU grace period完全覆盖。
在这种情况下，\co{r1}的最终值为1，\co{r2}的最终值为0。

然而，不可能出现\co{r1}的最终值为0且\co{r2}的最终值为1的情况。
这将意味着一个RCU read-side critical section完全覆盖了一个
grace period，这是被禁止的（或者至少构成RCU的一个bug）\@。
因此，RCU的等待读者保证有两个部分：
\begin{enumerate*}[(1)]
\item 如果给定RCU read-side critical section的任何部分先于
给定grace period的开始，则该critical section的全部先于该
grace period的结束。
\item 如果给定RCU read-side critical section的任何部分后于
给定grace period的结束，则该critical section的全部后于该
grace period的开始。
\end{enumerate*}
这个定义对几乎所有基于RCU的算法都是足够的，但对于想要更多的人，
简单的可执行RCU形式化模型作为Linux内核v4.17及更高版本的一部分
可供使用，如\cref{sec:formal:Axiomatic Approaches and RCU}中所述。
此外，RCU的排序属性在\cref{sec:memorder:RCU}中有更详细的讨论。

\QuickQuiz{
	如果\co{P0()}在
	\crefrange{fig:defer:RCU Reader and Later Grace Period}{fig:defer:RCU Reader Within Grace Period}
	中的访问是存储操作，会发生什么？
}\QuickQuizAnswer{
	完全相同的排序规则适用，即：
	(1)~如果\co{P0()}的RCU read-side critical section的任何部分
	先于\co{P1()}的grace period的开始，则\co{P0()}的
	RCU read-side critical section的全部将先于\co{P1()}的
	grace period的结束，
	(2)~如果\co{P0()}的RCU read-side critical section的任何部分
	后于\co{P1()}的grace period的结束，则\co{P0()}的
	RCU read-side critical section的全部将后于\co{P1()}的
	grace period的开始。

	在RCU read-side critical section中包含写操作可能看起来
	很奇怪，但这种能力不仅是被允许的，而且非常有用。
	例如，Linux内核经常执行RCU保护的链式数据结构遍历，
	然后获取目标数据元素的引用。
	因为该数据元素在此期间不能被释放，所以该元素的引用计数
	必须在遍历的RCU read-side critical section内递增。
	然而，该递增操作涉及对内存的写入。
	因此，在RCU read-side critical section中允许内存写入
	是一件非常好的事情。

	如果在RCU read-side critical section中进行写操作仍然
	看起来很奇怪，请回顾
	\cref{sec:count:Applying Exact Limit Counters}，
	其中展示了在读写锁读侧critical section中进行写操作的
	用例。
}\QuickQuizEnd

虽然RCU的等待读者能力确实有时被用于如
\crefrange{fig:defer:RCU Reader and Later Grace Period}
{fig:defer:RCU Reader Within Grace Period}
所示的变量赋值排序，但它更常被用于安全释放从链式结构中
移除的数据元素，正如在
\cref{sec:defer:Introduction to RCU}中所做的那样。
一般过程由以下伪代码说明：

\begin{enumerate}
\item	进行更改，例如，从链表中移除一个元素。
\item	等待所有已有的RCU read-side critical section完全完成
	（例如，通过使用\co{synchronize_rcu()}）。
\item	清理，例如，释放上面被替换的元素。
\end{enumerate}

\begin{figure}
\centering
\IfEbookSize{
\resizebox{2in}{!}{\includegraphics{defer/RCUGPorderingSummary}}
}{
\resizebox{2.5in}{!}{\includegraphics{defer/RCUGPorderingSummary}}
}
\caption{RCU Grace-Period排序保证总结}
\label{fig:defer:Summary of RCU Grace-Period Ordering Guarantees}
\end{figure}

这个更抽象的过程需要比
\crefrange{fig:defer:RCU Reader and Later Grace Period}
{fig:defer:RCU Reader Within Grace Period}
更抽象的图，因为那些图是针对特定石蕊测试的。
毕竟，无论RCU更新和RCU read-side critical section的形式如何，
RCU实现都必须正确工作。
\Cref{fig:defer:Summary of RCU Grace-Period Ordering Guarantees}
满足了这一需求，展示了四种可能的场景，每个场景中时间从上到下推进。
在每个场景中，RCU读者由左侧的方框栈表示，RCU更新者由右侧的
方框栈表示。

在第一个场景中，读者在更新者开始移除之前开始执行，因此该读者
可能持有对被移除数据元素的引用。
因此，更新者在读者完成之前不能释放此元素。
在第二个场景中，读者直到移除完成之后才开始执行。
该读者不可能获得对已移除数据元素的引用，因此可以在读者完成之前
释放该元素。
第三个场景类似于第二个，但说明了即使读者不可能获得对某元素的引用，
仍然允许将该元素的释放推迟到读者完成之后。
在第四个也是最后一个场景中，读者在更新者开始移除数据元素之前
开始执行，但该元素（错误地）在读者完成之前被释放。
正确的RCU实现不会允许第四种场景发生。
因此，此图说明了RCU的等待读者功能：
给定一个grace period，每个读者要么在该grace period结束之前结束，
要么在该grace period开始之后开始，或者两者兼有，
在这种情况下它完全包含在该grace period内。

因为RCU读者可以在更新进行的同时向前推进，不同的读者可能对
数据结构的状态持不同意见，这一主题将在下一节中讨论。

\subsubsection{维护最近更新对象的多个版本}
\label{sec:defer:Maintain Multiple Versions of Recently Updated Objects}

本节讨论RCU如何通过维护数据的多个版本来适应无同步读者。
因为这些无同步读者提供的时间同步非常弱，RCU用户通过空间同步
来补偿。
空间同步在\cref{chp:Partitioning and Synchronization Design}中讨论过，
并在实践中被大量使用以获得良好的性能和可扩展性。
在本节中，空间同步将被用来实现一种弱（但有用的）正确性形式，
以及出色的性能和可扩展性。

\cref{sec:defer:Minimal Insertion and Deletion}中的
\Cref{fig:defer:Deletion With Concurrent Readers}
展示了空间同步的一个简单变体，其中与\co{del_route()}
（见\cref{lst:defer:Insertion and Deletion With Concurrent Readers}）
并发运行的不同读者可能看到旧的\co{route}结构或空列表，
但无论哪种方式都能得到有效结果。
当然，仔细查看
\cref{fig:defer:Insertion With Concurrent Readers}
可以发现，\co{ins_route()}的调用也可能导致并发读者看到不同版本：
要么是初始的空列表，要么是新插入的\co{route}结构。
注意，reference counting
（\cref{sec:defer:Reference Counting}）
和hazard pointer
（\cref{sec:defer:Hazard Pointers}）
也可能导致并发读者看到不同版本，但RCU的轻量级读者使得这种情况
更为可能。

\begin{figure}
\centering
\resizebox{3in}{!}{\includegraphics{defer/multver}}
\caption{RCU数据结构的多个版本}
\label{fig:defer:Multiple RCU Data-Structure Versions}
\end{figure}

然而，维护多个弱一致性版本可能带来一些意外。
例如，考虑
\cref{fig:defer:Multiple RCU Data-Structure Versions}，
其中一个读者正在遍历一个被并发更新的链表。\footnote{
	RCU链表API可在
	\cref{sec:defer:RCU Linux-Kernel API}中找到。}
在图的第一行，读者引用数据项~A，
在第二行，它前进到~B，到目前为止看到的是A后面跟着~B\@。
在第三行，更新者移除了元素~A，在第四行，
更新者在列表末尾添加了元素~E。
在第五行也就是最后一行，读者完成了遍历，
看到的元素为~A到~E\@。

只是这样的列表从未在任何时刻存在过。
这种情况可能比
\cref{fig:defer:Deletion With Concurrent Readers}
所示的情况更令人惊讶，在那里不同的并发读者看到不同版本。
相比之下，在
\cref{fig:defer:Multiple RCU Data-Structure Versions}
中，读者看到的是一个从未实际存在过的版本！

解决这种奇怪情况的一种方法是通过较弱的语义。
读者遍历必须遇到在整个遍历期间都存在的任何数据项
（B、C和~D），但可能遇到也可能不遇到只在遍历的一部分时间
存在的数据项（A和~E）\@。
因此，在这种特定情况下，读者遍历遇到所有五个元素是完全合法的。
如果这个结果有问题，解决这种情况的另一种方法是使用更强的
同步机制，如读写锁，或巧妙地使用时间戳和版本控制，
如\cref{sec:defer:Quasi Multi-Version Concurrency Control}中所述。
当然，更强的机制会更昂贵，但工程生活就是关于选择和权衡的。

虽然这种情况看起来很奇怪，但它与现实世界完全一致。
正如我们在\cref{sec:cpu:Overheads}中看到的，
光速有限这一事实在计算机系统内不能被忽略，
在系统外部更不能被忽略。
这反过来意味着系统内表示系统外部真实世界状态的任何数据
始终是过时的，因此与真实世界不一致。
因此，序列\{A, B, C, D, E\}完全有可能在真实世界中发生过，
但由于光速延迟而从未在计算机系统的内存中体现。
在这种情况下，读者令人惊讶的遍历结果将正确反映现实。

因此，操作真实世界数据的算法必须考虑不一致的数据，
要么容忍不一致性，要么采取步骤排除或拒绝它们。
在许多情况下，这些算法也完全能够处理系统内部的不一致性。

\cref{sec:defer:Running Example}中介绍的pre-BSD数据包路由
示例就是一个很好的例证。
路由列表的内容由路由协议设置，这些协议具有显著的延迟
（秒甚至分钟级别），以避免路由不稳定。
因此，一旦路由更新到达给定系统，它可能已经将数据包发送到
错误的方向相当长时间了。
在更新传播的几微秒内再多发送几个错误方向的数据包显然不是问题，
因为处理延迟路由更新的相同高层协议操作也会处理内部不一致性。

互联网路由也不是唯一容忍不一致性的情况。
重复一下，任何系统中的数据跟踪系统外部状态的算法都必须
容忍不一致性，这包括安全策略（通常由人类委员会制定）、
存储配置和WiFi接入点，更不用说可移除硬件如麦克风、耳机、
摄像头、鼠标、打印机等等。
此外，\cref{fig:defer:RCU Usage in the Linux Kernel}所示的
大量Linux内核RCU API使用，
加上Linux内核大量使用reference counting以及其他项目越来越多地
使用\IXpl{hazard pointer}，表明对此类不一致性的容忍比人们
想象的更为常见。

这种常见情况下对不一致性的容忍的一个根本原因是，
单项查找在实践中比全数据结构遍历更为常见。
毕竟，全数据结构遍历比单项查找要昂贵得多，
因此开发者有动力避免此类遍历。
并发更新不仅比全遍历更不容易影响单项查找，
而且孤立的单项查找也无法检测到此类不一致性。
因此，在常见情况下，此类不一致性不仅是可容忍的，
实际上是不可见的。

在这种情况下，RCU读者可以被认为与更新者完全有序，
尽管这些读者可能正在执行与单线程程序完全相同的机器指令序列，
正如在\cpageref{sec:defer:Mysteries RCU}中所暗示的。
例如，回顾
\cpageref{lst:defer:Insertion and Deletion With Concurrent Readers}上的
\cref{lst:defer:Insertion and Deletion With Concurrent Readers}，
假设每个读者线程在其生命周期中恰好调用一次\co{access_route()}，
并且读者和更新者线程之间没有其他通信。
\begin{fcvref}[ln:defer:Insertion and Deletion With Concurrent Readers]
那么每次\co{access_route()}的调用可以排在产生该\co{route}结构
的\co{ins_route()}调用之后（该结构由清单中\co{access_route()}的
\clnref{access_rp}行访问），并排在任何后续的
\co{ins_route()}或\co{del_route()}调用之前。
\end{fcvref}

总之，维护多个版本正是使RCU读者具有极低开销的关键，
而且如前所述，许多算法不受多版本的影响。
然而，确实有些算法完全无法处理多个版本。
有一些技术可以使此类算法适应
RCU~\cite{PaulEdwardMcKenneyPhD}，例如
\cref{sec:together:Correlated Data Elements}中描述的sequence locking的
使用。

\paragraph{练习}
\label{sec:defer:Exercises}

这些示例假设在整个更新操作期间持有互斥锁（exclusive lock），
这意味着在给定时间最多有两个版本的列表处于活动状态。

\QuickQuizSeries{%
\QuickQuizB{
	你将如何修改删除示例以允许列表有两个以上的
	活动版本？
}\QuickQuizAnswerB{
	实现这一目标的一种方式如
	\cref{lst:defer:Concurrent RCU Deletion}所示。

\begin{listing}
\begin{VerbatimL}
spin_lock(&mylock);
p = search(head, key);
if (p == NULL)
	spin_unlock(&mylock);
else {
	list_del_rcu(&p->list);
	spin_unlock(&mylock);
	synchronize_rcu();
	kfree(p);
}
\end{VerbatimL}
\caption{并发RCU删除}
\label{lst:defer:Concurrent RCU Deletion}
\end{listing}

	注意，这意味着多个并发删除可能在
	\co{synchronize_rcu()}中等待。
}\QuickQuizEndB
%
\QuickQuizM{
	在任何给定时间，给定列表可以有多少个RCU版本
	处于活动状态？
}\QuickQuizAnswerM{
	这取决于同步设计。
	如果保护更新的信号量在grace period期间一直持有，
	则最多可以有两个版本——旧版本和新版本。

	然而，假设只有搜索、更新和\co{list_replace_rcu()}
	受锁保护，使得\co{synchronize_rcu()}在该锁之外，
	类似于\cref{lst:defer:Concurrent RCU Deletion}所示的代码。
	进一步假设大量线程大约在同一时间进行RCU替换，
	并且读者也在不断遍历数据结构。

	那么以下事件序列可能发生，从
	\cref{fig:defer:Multiple RCU Data-Structure Versions}
	的结束状态开始：

	\begin{enumerate}
	\item	线程~A遍历列表，获取对元素~C的引用。
	\item	线程~B将元素~C替换为新元素~F，
		然后等待其\co{synchronize_rcu()}调用返回。
	\item	线程~C遍历列表，获取对元素~F的引用。
	\item	线程~D将元素~F替换为新元素~G，
		然后等待其\co{synchronize_rcu()}调用返回。
	\item	线程~E遍历列表，获取对元素~G的引用。
	\item	线程~F将元素~G替换为新元素~H，
		然后等待其\co{synchronize_rcu()}调用返回。
	\item	线程~G遍历列表，获取对元素~H的引用。
	\item	前两步用新的元素快速重复，使得所有这些
		操作在任何\co{synchronize_rcu()}调用返回之前完成。
	\end{enumerate}

	因此，可以有任意数量的活动版本，仅受内存和在一个
	grace period内可以完成多少次更新的限制。
	但请注意，更新如此频繁的数据结构不太可能是RCU的
	好候选者\@。
	尽管如此，RCU在必要时可以处理高更新率。
}\QuickQuizEndM
%
\QuickQuizE{
	如何降低RCU每次更新的开销？
}\QuickQuizAnswerE{
	降低RCU每次更新开销最有效的方法是增加由
	给定grace period服务的更新数量。
	这之所以有效，是因为每个grace period的开销几乎
	与该grace period服务的更新数量无关。

	一种方法是延迟给定grace period的开始，
	希望在此期间有更多需要该grace period的更新出现。
	另一种方法是减慢grace period的执行速度，
	希望在此期间有更多需要额外grace period的更新积累。

	还有许多其他可能的优化，狂热的读者可以参考
	Linux内核的RCU实现。
}\QuickQuizEndE
}

\subsubsection{RCU基础总结}
\label{sec:defer:Summary of RCU Fundamentals}

本节描述了基于RCU的链式结构算法的三个基本组件：

\begin{enumerate}
\item	用于添加新数据的publish-subscribe机制，
	更新侧使用\co{rcu_assign_pointer()}进行发布，
	读侧使用\co{rcu_dereference()}进行订阅，

\item	一种等待已有RCU读者完成的方式，
	读者一侧由\co{rcu_read_lock()}和\co{rcu_read_unlock()}界定，
	更新者一侧通过\co{synchronize_rcu()}或
	\co{call_rcu()}等待
	（形式化描述见\cref{sec:memorder:RCU}），
	以及

\item	一种维护多版本的规范，以允许更改而不伤害或
	过度延迟并发RCU读者。
\end{enumerate}

\QuickQuiz{
	既然\co{rcu_read_lock()}和\co{rcu_read_unlock()}
	都不会自旋或阻塞，RCU更新者怎么可能延迟RCU读者？
}\QuickQuizAnswer{
	给定RCU更新者所进行的修改将导致相应CPU使包含数据的
	缓存行无效，迫使运行并发RCU读者的CPU产生昂贵的
	缓存未命中。
	（你能设计一种更改数据结构而\emph{不}对并发读者造成
	昂贵缓存未命中的算法吗？对后续读者呢？）
}\QuickQuizEnd

这三个RCU组件允许在并发读者可能正在执行与单线程实现中
读者完全相同的机器指令序列的情况下进行数据更新。
这些RCU组件可以以不同方式组合来实现种类繁多的不同类型的
基于RCU的算法，其中一些在
\cref{sec:defer:RCU Usage}中介绍。
然而，通常更好的做法是在更高的抽象层次上工作。
为此，下一节描述了Linux内核API，其中包括列表等简单数据结构。
